\documentclass[12pt,a4paper]{article}
\usepackage[utf8]{inputenc}
\usepackage[T1]{fontenc}
\usepackage{amsmath,amssymb,amsthm}
\usepackage{geometry}
\usepackage{enumitem}
\usepackage{hyperref}
\usepackage{fancyhdr}
\usepackage{titlesec}

\geometry{margin=1in}
\pagestyle{fancy}
\fancyhf{}
\rhead{Lecture 3}
\lhead{Quantum Systems via Linear Algebra}
\cfoot{\thepage}

\titleformat{\section}{\large\bfseries}{}{0em}{}
\titleformat{\subsection}{\normalsize\bfseries}{}{0em}{}

\newtheorem{theorem}{Theorem}
\newtheorem{definition}{Definition}

\title{\textbf{Lecture 3: Linear Operators and Self-Adjoint Matrices} \\ \large Matrix Elements, Eigendecomposition, and the $X$, $Y$, $Z$ Matrices}
\author{Quantum Systems via Linear Algebra, Differential Equations, and Probability}
\date{}

\begin{document}
\maketitle
\thispagestyle{fancy}

% ============================================================
\section{1. Expanding Vectors in a Basis}
% ============================================================

Let $\{\mathbf{e}_i\}$ be an \textbf{orthonormal basis} ($\mathbf{e}_j^\dagger\mathbf{e}_i=\delta_{ji}$) of dimension $N$.  Any vector can be expanded as
\[
    \mathbf{v} = \sum_{i=1}^{N} \alpha_i \,\mathbf{e}_i, \qquad \alpha_i = \mathbf{e}_i^\dagger\mathbf{v}.
\]
The second equation follows by left-multiplying by $\mathbf{e}_j^\dagger$: orthonormality collapses the sum to a single term $\alpha_j = \mathbf{e}_j^\dagger\mathbf{v}$.  Substituting back:
\[
    \boxed{\mathbf{v} = \sum_i \mathbf{e}_i\,(\mathbf{e}_i^\dagger\mathbf{v}).}
\]
The identity $\sum_i \mathbf{e}_i\mathbf{e}_i^\dagger = I$ is the \textbf{resolution of identity}---each term $\mathbf{e}_i\mathbf{e}_i^\dagger$ is a rank-1 projector onto the $i$-th basis direction.

% ============================================================
\section{2. Linear Operators}
% ============================================================

A linear operator maps column vectors to column vectors.  Given any input $\mathbf{v}$, it produces a unique output $\mathbf{w}=M\mathbf{v}$.

\begin{definition}[Linear operator]
A \textbf{linear operator} $M$ is a map from vectors to vectors satisfying two linearity rules:
\begin{align}
    M(z\mathbf{v}) &= z\,M\mathbf{v}, \quad \forall\, z\in\mathbb{C}, \\
    M(\mathbf{u}+\mathbf{v}) &= M\mathbf{u} + M\mathbf{v}.
\end{align}
In words: complex constants can be pulled through the operator, and the operator distributes over sums.
\end{definition}

\subsection{2.1 Matrix elements}
Given a basis $\{\mathbf{e}_i\}$, the operator $M$ is fully characterized by its \textbf{matrix elements}:
\[
    M_{ij} \equiv \mathbf{e}_i^\dagger\,M\,\mathbf{e}_j.
\]
This is a concrete complex number obtained by: (1) applying $M$ to basis vector $\mathbf{e}_j$ to get a new vector, then (2) taking the inner product with $\mathbf{e}_i^\dagger$.

If $M\mathbf{v}=\mathbf{w}$, insert the resolution of identity (from \S1: $I = \sum_j \mathbf{e}_j\mathbf{e}_j^\dagger$) between $M$ and $\mathbf{v}$:
\[
    w_i = \mathbf{e}_i^\dagger\mathbf{w} = \mathbf{e}_i^\dagger M\mathbf{v}
    = \mathbf{e}_i^\dagger M\!\underbrace{\Bigl(\sum_j \mathbf{e}_j\mathbf{e}_j^\dagger\Bigr)}_{=\,I}\!\mathbf{v}
    = \sum_j (\mathbf{e}_i^\dagger M\mathbf{e}_j)(\mathbf{e}_j^\dagger\mathbf{v}) = \sum_j M_{ij}\,v_j,
\]
which is simply matrix--column-vector multiplication:
\[
    \begin{pmatrix}w_1\\w_2\\\vdots\end{pmatrix}
    = \begin{pmatrix}M_{11}&M_{12}&\cdots\\M_{21}&M_{22}&\cdots\\\vdots&\vdots&\ddots\end{pmatrix}
    \begin{pmatrix}v_1\\v_2\\\vdots\end{pmatrix}.
\]

\medskip
\noindent\textbf{Basis dependence:} The specific numbers $M_{ij}$ depend on the choice of basis.  A different basis gives different matrix elements for the same abstract operator.  This is analogous to how the components of a vector change when you rotate your coordinate axes.

\medskip
\noindent\textbf{What doesn't change:} The \emph{eigenvalues} of $M$ are intrinsic to the operator---they are the same in every basis.  (So are the trace $\operatorname{Tr}(M) = \sum_i M_{ii}$, the determinant, and the rank.)  This is why measurement outcomes, which are eigenvalues by Axiom~2, are physically meaningful: they are basis-independent facts about the operator, not coordinate-dependent artifacts.

\subsection{2.2 Quadratic form}
For any vectors $\mathbf{u}$, $\mathbf{v}$ and operator $M$, the scalar $\mathbf{u}^\dagger M\mathbf{v}$ is unambiguous---it reduces to a sum over matrix elements.  The special case $\mathbf{v}^\dagger M\mathbf{v}$ is the \textbf{Rayleigh quotient}, already familiar from linear algebra.

% ============================================================
\section{3. Hermitian Conjugate (Adjoint)}
% ============================================================

\begin{definition}[Hermitian conjugate]
For every matrix $M$, there is a unique matrix $M^\dagger$ (the \textbf{Hermitian conjugate}, or \textbf{adjoint}) satisfying: for all $\mathbf{u},\mathbf{v}$,
\[
    \mathbf{u}^\dagger(M\mathbf{v}) = (M^\dagger\mathbf{u})^\dagger\mathbf{v}.
\]
\end{definition}

\noindent The adjoint is the operator that, when moved from the right side of an inner product to the left side, takes $M$'s place.  For numbers, the analog is $u^*(zv) = (z^* u)^* v$.

\medskip
\noindent The dagger operation is the matrix analog of complex conjugation.  The full analogy:
\[
\renewcommand{\arraystretch}{1.3}
\begin{array}{c|c}
\textbf{Numbers} & \textbf{Matrices} \\\hline
z^* & M^\dagger \\
z = z^* \;(\text{real}) & M = M^\dagger \;(\text{self-adjoint}) \\
|z|^2 = z^* z & \|\mathbf{v}\|^2 = \mathbf{v}^\dagger\mathbf{v} \\
|z| = 1 \;(\text{unit complex}) & U^\dagger U = I \;(\text{unitary})
\end{array}
\]
The last row---unit complex numbers $\leftrightarrow$ unitary matrices---will do real work when we reach time evolution.

\subsection{3.1 Matrix elements of the adjoint}
\[
    \boxed{(M^\dagger)_{ij} = M_{ji}^*.}
\]
Recipe: \textbf{transpose} (swap $i\leftrightarrow j$) and \textbf{complex conjugate}.  The transpose interchanges rows and columns; then you conjugate every entry.

\subsection{3.2 Example}
\[
    M = \begin{pmatrix}2 & 6+i\\7 & 4-i\end{pmatrix}
    \;\xrightarrow{\text{transpose}}\;
    \begin{pmatrix}2 & 7\\6+i & 4-i\end{pmatrix}
    \;\xrightarrow{\text{conjugate}}\;
    M^\dagger = \begin{pmatrix}2 & 7\\6-i & 4+i\end{pmatrix}.
\]

% ============================================================
\section{4. Self-Adjoint Matrices}
% ============================================================

\begin{definition}[Self-adjoint matrix]
A matrix $M$ is \textbf{self-adjoint} (Hermitian) if
\[
    M = M^\dagger \quad\Longleftrightarrow\quad M_{ij} = M_{ji}^*.
\]
\end{definition}

Self-adjoint is the matrix analogue of ``real'' for numbers ($z=z^*$).  Properties:
\begin{itemize}[nosep]
    \item \textbf{Diagonal elements are real}: $M_{ii} = M_{ii}^*$, so $M_{ii}\in\mathbb{R}$.
    \item \textbf{Off-diagonal elements} come in conjugate pairs: $M_{ij} = M_{ji}^*$.
\end{itemize}

\noindent\textbf{Example:}
\[
    \begin{pmatrix}1 & i & 4-i\\ -i & 7 & 5\\ 4+i & 5 & -2.6\end{pmatrix}
\]
is self-adjoint: diagonal entries ($1,7,-2.6$) are real; $(1,2)$-entry $i$ and $(2,1)$-entry $-i$ are conjugates; etc.

\medskip
\noindent\textbf{Central role (Axiom~2):} Self-adjoint matrices represent observables.  The result of any measurement is a real number, and self-adjoint matrices guarantee real eigenvalues.

% ============================================================
\section{5. Eigenvalues and Eigenvectors}
% ============================================================

\begin{definition}[Eigenvector, eigenvalue]
A nonzero vector $\mathbf{e}_\lambda$ is an \textbf{eigenvector} of $M$ with \textbf{eigenvalue} $\lambda$ if
\[
    M\mathbf{e}_\lambda = \lambda\,\mathbf{e}_\lambda.
\]
The matrix merely \emph{rescales} the vector without changing its direction.
\end{definition}

\noindent\textbf{Note:} Eigenvectors are defined only up to a multiplicative constant.  If $\mathbf{e}_\lambda$ is an eigenvector, so is $z\mathbf{e}_\lambda$ for any $z\neq 0$, with the same eigenvalue.

\subsection{5.1 Spectral theorem for self-adjoint matrices}

\begin{theorem}
If $M = M^\dagger$ on an $N$-dimensional space, then:
\begin{enumerate}[nosep]
    \item All eigenvalues $\lambda_i$ are \textbf{real}.
    \item Eigenvectors with \textbf{distinct} eigenvalues are \textbf{orthogonal}.
    \item There exists a \textbf{complete} orthonormal basis of eigenvectors (i.e., $N$ of them).
\end{enumerate}
\end{theorem}

\noindent This is the most important theorem for our purposes.  It guarantees that the eigenvectors of any observable form a basis, and that measurement outcomes (eigenvalues) are always real.

\subsection{5.2 Proof: eigenvalues are real}
Start from $M\mathbf{e}_i=\lambda_i\mathbf{e}_i$.  Take the conjugate transpose:
\[
    \mathbf{e}_i^\dagger M^\dagger = \lambda_i^*\mathbf{e}_i^\dagger.
\]
Since $M=M^\dagger$, both equations involve the same matrix $M$.  Left-multiply the first by $\mathbf{e}_i^\dagger$ and right-multiply the second by $\mathbf{e}_i$:
\[
    \mathbf{e}_i^\dagger M\mathbf{e}_i = \lambda_i\,\mathbf{e}_i^\dagger\mathbf{e}_i, \qquad
    \mathbf{e}_i^\dagger M\mathbf{e}_i = \lambda_i^*\,\mathbf{e}_i^\dagger\mathbf{e}_i.
\]
Cancelling $\mathbf{e}_i^\dagger\mathbf{e}_i\neq 0$ gives $\lambda_i = \lambda_i^*$, so $\lambda_i\in\mathbb{R}$.\;$\square$

\subsection{5.3 Proof: eigenvectors with distinct eigenvalues are orthogonal}
Let $M\mathbf{e}_i=\lambda_i\mathbf{e}_i$ and $M\mathbf{e}_j=\lambda_j\mathbf{e}_j$ with $\lambda_i\neq\lambda_j$.  The conjugate transpose of the second:
\[
    \mathbf{e}_j^\dagger M = \lambda_j\,\mathbf{e}_j^\dagger \quad\text{(using $\lambda_j^*=\lambda_j$)}.
\]
Compute $\mathbf{e}_j^\dagger M\mathbf{e}_i$ two ways:
\[
    \mathbf{e}_j^\dagger M\mathbf{e}_i = \lambda_i\,\mathbf{e}_j^\dagger\mathbf{e}_i, \qquad
    \mathbf{e}_j^\dagger M\mathbf{e}_i = \lambda_j\,\mathbf{e}_j^\dagger\mathbf{e}_i.
\]
Subtracting: $(\lambda_i - \lambda_j)\,\mathbf{e}_j^\dagger\mathbf{e}_i=0$.  Since $\lambda_i\neq\lambda_j$, we conclude $\mathbf{e}_j^\dagger\mathbf{e}_i=0$.\;$\square$

\medskip
\noindent\textbf{Physical meaning:} Orthogonal eigenvectors correspond to distinguishable states---states that can be told apart with certainty by a single measurement.

% ============================================================
\section{6. Connection to the Axioms}
% ============================================================

With the spectral theorem in hand, the axioms from Lecture~1 become concrete:

\begin{enumerate}[nosep]
    \item \textbf{Axiom~2 (Observable):} Every measurable quantity is a self-adjoint matrix $M = M^\dagger$.  Its eigenvalues $\lambda_i \in \mathbb{R}$ are the only possible outcomes.
    \item \textbf{Eigenvectors are states of definite value.}  The eigenvector $\mathbf{e}_i$ corresponding to $\lambda_i$ is the state in which $M$ has the unambiguous value $\lambda_i$.
    \item \textbf{Axiom~3 (Born rule):} If the system is in state $\mathbf{v}$ and we measure $M$, the probability of obtaining $\lambda_i$ is
    \[
        P(\lambda_i) = |\mathbf{e}_i^\dagger\mathbf{v}|^2.
    \]
    The complex number $\mathbf{e}_i^\dagger\mathbf{v}$ is the \textbf{probability amplitude}; squaring its magnitude gives the probability.
\end{enumerate}

\noindent\textbf{Normalization.}  Since probabilities must sum to 1:
\[
    \sum_i P(\lambda_i) = \sum_i |\mathbf{e}_i^\dagger\mathbf{v}|^2 = \mathbf{v}^\dagger\mathbf{v} = 1.
\]
Physical states are unit vectors.  The overall phase $e^{i\theta}\mathbf{v}$ is physically equivalent to $\mathbf{v}$ (all probabilities are unchanged).

% ============================================================
\section{7. The $X$, $Y$, $Z$ Matrices}
% ============================================================

The three standard observables for a two-level system are self-adjoint matrices on $\mathbb{C}^2$.  We work in the $\{\mathbf{e}_0,\mathbf{e}_1\}$ basis and compute each matrix from its known eigenvectors and eigenvalues (derived in Lecture~2).

\subsection{7.1 Computing $Z$}
We know: $Z\mathbf{e}_0=+1\cdot\mathbf{e}_0$ and $Z\mathbf{e}_1=-1\cdot\mathbf{e}_1$.  The matrix elements:
\begin{align}
    Z_{00} &= \mathbf{e}_0^\dagger Z\mathbf{e}_0 = +1\cdot\mathbf{e}_0^\dagger\mathbf{e}_0 = +1, \\
    Z_{01} &= \mathbf{e}_0^\dagger Z\mathbf{e}_1 = -1\cdot\mathbf{e}_0^\dagger\mathbf{e}_1 = 0, \\
    Z_{10} &= \mathbf{e}_1^\dagger Z\mathbf{e}_0 = +1\cdot\mathbf{e}_1^\dagger\mathbf{e}_0 = 0, \\
    Z_{11} &= \mathbf{e}_1^\dagger Z\mathbf{e}_1 = -1\cdot\mathbf{e}_1^\dagger\mathbf{e}_1 = -1.
\end{align}
Therefore:
\[
    Z = \begin{pmatrix}1&0\\0&-1\end{pmatrix}.
\]

\subsection{7.2 Computing $X$}
The eigenvectors of $X$ are $\mathbf{x}_+=\frac{1}{\sqrt{2}}(\mathbf{e}_0+\mathbf{e}_1)$ with eigenvalue $+1$ and $\mathbf{x}_-=\frac{1}{\sqrt{2}}(\mathbf{e}_0-\mathbf{e}_1)$ with eigenvalue $-1$.  Express $\mathbf{e}_0$ and $\mathbf{e}_1$ in terms of $\mathbf{x}_\pm$:
\[
    \mathbf{e}_0 = \frac{1}{\sqrt{2}}(\mathbf{x}_++\mathbf{x}_-), \qquad
    \mathbf{e}_1 = \frac{1}{\sqrt{2}}(\mathbf{x}_+-\mathbf{x}_-).
\]
The cleanest approach is to apply $X$ to each basis vector first, then take inner products.  Since $X\mathbf{x}_+=\mathbf{x}_+$ and $X\mathbf{x}_-=-\mathbf{x}_-$:
\[
    X\mathbf{e}_0 = \frac{1}{\sqrt{2}}(X\mathbf{x}_++X\mathbf{x}_-) = \frac{1}{\sqrt{2}}(\mathbf{x}_+-\mathbf{x}_-) = \mathbf{e}_1,
\]
\[
    X\mathbf{e}_1 = \frac{1}{\sqrt{2}}(X\mathbf{x}_+-X\mathbf{x}_-) = \frac{1}{\sqrt{2}}(\mathbf{x}_++\mathbf{x}_-) = \mathbf{e}_0.
\]
Now the matrix elements follow immediately:
\begin{align}
    X_{00} &= \mathbf{e}_0^\dagger(X\mathbf{e}_0) = \mathbf{e}_0^\dagger\mathbf{e}_1 = 0, &\quad
    X_{01} &= \mathbf{e}_0^\dagger(X\mathbf{e}_1) = \mathbf{e}_0^\dagger\mathbf{e}_0 = 1, \\[4pt]
    X_{10} &= \mathbf{e}_1^\dagger(X\mathbf{e}_0) = \mathbf{e}_1^\dagger\mathbf{e}_1 = 1, &\quad
    X_{11} &= \mathbf{e}_1^\dagger(X\mathbf{e}_1) = \mathbf{e}_1^\dagger\mathbf{e}_0 = 0.
\end{align}
Therefore:
\[
    X = \begin{pmatrix}0&1\\1&0\end{pmatrix}.
\]

\subsection{7.3 Computing $Y$}
The eigenvectors are $\mathbf{y}_+=\frac{1}{\sqrt{2}}(\mathbf{e}_0+i\mathbf{e}_1)$ with eigenvalue $+1$ and $\mathbf{y}_-=\frac{1}{\sqrt{2}}(\mathbf{e}_0-i\mathbf{e}_1)$ with eigenvalue $-1$.  Express the basis vectors:
\[
    \mathbf{e}_0 = \frac{1}{\sqrt{2}}(\mathbf{y}_++\mathbf{y}_-), \qquad
    \mathbf{e}_1 = \frac{1}{i\sqrt{2}}(\mathbf{y}_+-\mathbf{y}_-) = \frac{-i}{\sqrt{2}}(\mathbf{y}_+-\mathbf{y}_-).
\]
Apply $Y$ to each basis vector:
\[
    Y\mathbf{e}_0 = \frac{1}{\sqrt{2}}(Y\mathbf{y}_++Y\mathbf{y}_-) = \frac{1}{\sqrt{2}}(\mathbf{y}_+-\mathbf{y}_-) = i\,\mathbf{e}_1,
\]
\[
    Y\mathbf{e}_1 = \frac{-i}{\sqrt{2}}(Y\mathbf{y}_+-Y\mathbf{y}_-) = \frac{-i}{\sqrt{2}}(\mathbf{y}_++\mathbf{y}_-) = -i\,\mathbf{e}_0.
\]
The matrix elements:
\begin{align}
    Y_{00} &= \mathbf{e}_0^\dagger(Y\mathbf{e}_0) = i\,\mathbf{e}_0^\dagger\mathbf{e}_1 = 0, &\quad
    Y_{01} &= \mathbf{e}_0^\dagger(Y\mathbf{e}_1) = -i\,\mathbf{e}_0^\dagger\mathbf{e}_0 = -i, \\[4pt]
    Y_{10} &= \mathbf{e}_1^\dagger(Y\mathbf{e}_0) = i\,\mathbf{e}_1^\dagger\mathbf{e}_1 = i, &\quad
    Y_{11} &= \mathbf{e}_1^\dagger(Y\mathbf{e}_1) = -i\,\mathbf{e}_1^\dagger\mathbf{e}_0 = 0.
\end{align}
Therefore:
\[
    Y = \begin{pmatrix}0&-i\\i&0\end{pmatrix}.
\]
The imaginary entries arise because the $Y$ eigenvectors have factors of $i$---this is the matrix-level consequence of the argument from Lecture~2, \S3.3 that complex coefficients are essential for the third basis.

\subsection{7.4 Summary of the $X$, $Y$, $Z$ matrices}
\[
    \boxed{Z = \begin{pmatrix}1&0\\0&-1\end{pmatrix}, \qquad
    X = \begin{pmatrix}0&1\\1&0\end{pmatrix}, \qquad
    Y = \begin{pmatrix}0&-i\\i&0\end{pmatrix}.}
\]
Each is self-adjoint ($M = M^\dagger$), each has eigenvalues $+1$ and $-1$, and each has orthogonal eigenvectors forming a complete basis.  (In physics these are called the Pauli matrices and denoted $\sigma_z, \sigma_x, \sigma_y$.)

\subsection{7.5 Eigenvalues and eigenvectors summary}
\begin{align}
    Z: &\quad \mathbf{e}_0\;(\lambda=+1),\quad \mathbf{e}_1\;(\lambda=-1). \\
    X: &\quad \mathbf{x}_+=\tfrac{1}{\sqrt{2}}(\mathbf{e}_0+\mathbf{e}_1)\;(\lambda=+1),\quad \mathbf{x}_-=\tfrac{1}{\sqrt{2}}(\mathbf{e}_0-\mathbf{e}_1)\;(\lambda=-1). \\
    Y: &\quad \mathbf{y}_+=\tfrac{1}{\sqrt{2}}(\mathbf{e}_0+i\mathbf{e}_1)\;(\lambda=+1),\quad \mathbf{y}_-=\tfrac{1}{\sqrt{2}}(\mathbf{e}_0-i\mathbf{e}_1)\;(\lambda=-1).
\end{align}

\subsection{7.6 Algebraic properties (with derivations)}

\noindent\textbf{Each Pauli matrix squares to the identity.}  For $X$:
\[
    X^2 = \begin{pmatrix}0&1\\1&0\end{pmatrix}\begin{pmatrix}0&1\\1&0\end{pmatrix}
    = \begin{pmatrix}1&0\\0&1\end{pmatrix} = I.\;\checkmark
\]
Similarly $Y^2 = Z^2 = I$.  (Equivalently: each Pauli matrix has eigenvalues $\pm 1$, so its square has eigenvalues $(\pm 1)^2 = 1$, which means $M^2 = I$.)

\medskip
\noindent\textbf{Products of distinct Pauli matrices.}  Compute $XY$ explicitly:
\[
    XY = \begin{pmatrix}0&1\\1&0\end{pmatrix}\begin{pmatrix}0&-i\\i&0\end{pmatrix}
    = \begin{pmatrix}i&0\\0&-i\end{pmatrix}
    = i\begin{pmatrix}1&0\\0&-1\end{pmatrix}
    = iZ.
\]
Since $X$ and $Y$ do not commute, the reversed product differs:
\[
    YX = \begin{pmatrix}0&-i\\i&0\end{pmatrix}\begin{pmatrix}0&1\\1&0\end{pmatrix}
    = \begin{pmatrix}-i&0\\0&i\end{pmatrix}
    = -iZ.
\]
Subtracting: $[X,Y] = XY - YX = iZ - (-iZ) = 2iZ$.  The full set of cyclic commutation relations:
\[
    \boxed{[X,Y] = 2iZ, \qquad [Y,Z] = 2iX, \qquad [Z,X] = 2iY.}
\]
These three relations are the defining relations of the \textbf{Lie algebra $\mathfrak{su}(2)$}---the algebraic structure underlying all spin and angular momentum in quantum mechanics.  Non-commutativity ($[M,N]\neq 0$) is the mathematical expression of the fact that you cannot simultaneously know the values of two different observables.

% ============================================================
\section{8. Summary}
% ============================================================

\begin{enumerate}[nosep]
    \item Any vector can be expanded in an orthonormal basis: $\mathbf{v}=\sum_i\mathbf{e}_i\,(\mathbf{e}_i^\dagger\mathbf{v})$.  The resolution of identity $\sum_i\mathbf{e}_i\mathbf{e}_i^\dagger=I$ is an indispensable computational tool.
    \item \textbf{Linear operators} map vectors to vectors; they are represented by matrices with elements $M_{ij}=\mathbf{e}_i^\dagger M\mathbf{e}_j$.
    \item The \textbf{adjoint} $M^\dagger$ is obtained by transposing and complex conjugating the matrix.
    \item \textbf{Self-adjoint} matrices ($M=M^\dagger$) have real eigenvalues and orthogonal eigenvectors---they represent observables (Axiom~2).
    \item The \textbf{axioms} in action: observables $\leftrightarrow$ self-adjoint matrices; outcomes $\leftrightarrow$ eigenvalues; probability of outcome $\lambda_i$ is $|\mathbf{e}_i^\dagger\mathbf{v}|^2$ (Axiom~3).
    \item The $X$, $Y$, $Z$ matrices are the self-adjoint matrices for the three standard observables of a two-level system, each with eigenvalues $\pm 1$, derived from the known eigenvectors.
\end{enumerate}

\end{document}
